\documentclass[11pt]{article}

\usepackage[margin=1in]{geometry}
\usepackage{hyperref}

\begin{document}

\textit{Author's note: AI was used to check grammar and rephrase for better wording. The prompt used was: ``Please check grammar, cohesiveness and correct wording of the following.''}

\vspace{1em}

\noindent\textbf{Mission Statement:}

A coordinated multi-agent system deployed to Europa's surface to autonomously characterize 
ice shell thickness and identify an optimal cryobot insertion site without real-time Earth
support. The destination is Europa, chosen because its subsurface ocean makes it one of the
most promising locations in the Solar System for astrobiology and future exploration. 
The mission will deploy a surface rover equipped with a ground-penetrating radar sounder, 
supported by two orbital reconnaissance agents that provide terrain mapping, communications 
relay, and cooperative localization. By leveraging radar sounding techniques analogous to the 
Europa Clipper REASON instrument, the system will generate detailed subsurface ice maps to 
support future cryobot missions.

\vspace{1em}

\noindent\textbf{Operational Phases:}

The mission consists of six main phases: Launch and Commissioning, Cruise, Europa Orbit 
Insertion and Reconnaissance, Landing, Surface Checkout and Site Characterization, Surface 
Survey and Ice Characterization, and End-of-Mission Operations. Launch and cruise prepare the
system for arrival at Europa. Upon arrival, the orbital agents conduct reconnaissance and 
identify a suitable landing region. Landing is followed by rover deployment, system checkout, 
and sensor calibration. Main surface operations focus on autonomous rover traversal while 
collecting radar measurements of the ice shell and building a subsurface map. The orbital 
agents continuously update terrain maps and relay communications. The final phase includes 
science data downlink, candidate cryobot site selection, optional extended operations, and 
safe system passivation.

\vspace{1em}

\noindent\textbf{Systems Roles:}

The rover operates autonomously during critical events such as landing, navigation, radar 
sounding, and site-selection activities. The two orbital agents provide terrain reconnaissance,
communications relay, cooperative localization support, and validation of shared mapping data. 
Ground teams remain in the loop during cruise, strategic mission planning, and anomaly 
resolution, but all tactical decisions must be performed onboard due to the approximately 
43--44 minute one-way communication delay between Earth and Jupiter. Radiation exposure, 
cryogenic temperatures, communication delays, and limited bandwidth are key operational
constraints. The radar sounder provides continuous subsurface measurements that support 
ice-thickness estimation and cryobot site identification throughout the mission.

\vspace{1em}

\noindent\textbf{References}

\begin{thebibliography}{9}

\bibitem{EuropaClipper}
NASA, ``Europa Clipper Mission Overview,'' NASA Science. [Online]. Available: \url{https://science.nasa.gov/mission/europa-clipper/}. [Accessed: Jun. 21, 2026].

\bibitem{REASON}
NASA, ``Radar for Europa Assessment and Sounding: Ocean to Near-surface (REASON),'' Europa Clipper Mission. [Online]. Available: \url{https://europa.nasa.gov/spacecraft/instruments/reason/}. [Accessed: Jun. 21, 2026].

\bibitem{DecadalSurvey}
National Academies of Sciences, Engineering, and Medicine, \textit{Origins, Worlds, and Life: A Decadal Strategy for Planetary Science and Astrobiology 2023--2032}. Washington, DC: National Academies Press, 2022.

\end{thebibliography}

\end{document}
```
